%% ===============================
%% anlagen/listings.tex
%% Quellcode-Auszüge: Multimodaler Sprach-Assistent
%% ===============================

% WICHTIG:
% KEIN \chapter hier! Wird in der Hauptdatei definiert.

% ===============================
% Listings Setup für Python
% ===============================

\lstdefinelanguage{Python}{
    morekeywords={as,assert,break,class,continue,def,del,elif,else,except,False,finally,for,from,global,if,import,in,is,lambda,None,nonlocal,not,or,pass,raise,return,True,try,while,with,yield},
    sensitive=true,
    morecomment=[l]{\#},
    morestring=[b]',
    morestring=[b]"
}

\lstset{
    language=Python,
    basicstyle=\ttfamily\small,
    numbers=left,
    numberstyle=\tiny\color{gray},
    stepnumber=1,
    numbersep=5pt,
    frame=single,
    rulecolor=\color{black!30},
    backgroundcolor=\color{white!97!gray},
    breaklines=true,
    captionpos=b,
    commentstyle=\color{green!50!black},
    keywordstyle=\color{blue},
    stringstyle=\color{orange}
}

% ===============================
% Speech-to-Text (Whisper)
% ===============================
\section{Speech-to-Text (Whisper)}

Die folgende Implementierung zeigt die erste Stufe der Pipeline, in der gesprochene Sprache mithilfe des Whisper-Modells in Text transformiert wird.

\begin{lstlisting}[caption={speech_to_text.py – Audio-Transkription mittels Whisper}]
import whisper
import tempfile

model = whisper.load_model("base")

def transcribe_audio(uploaded_file):
    with tempfile.NamedTemporaryFile(delete=False, suffix=".wav") as tmp:
        tmp.write(uploaded_file.read())
        tmp_path = tmp.name

    result = model.transcribe(tmp_path)
    return result["text"]
\end{lstlisting}

% ===============================
% NLP Analyse (spaCy)
% ===============================
\section{Natural Language Processing}

In diesem Schritt erfolgt die linguistische Analyse des transkribierten Textes mithilfe von spaCy. Der Text wird tokenisiert und grammatikalisch annotiert.

\begin{lstlisting}[caption={nlp_analysis.py – Tokenisierung und Part-of-Speech Tagging}]
import spacy

nlp = spacy.load("de_core_news_sm")

def analyze(text):
    doc = nlp(text)

    tokens = [{"text": t.text, "pos": t.pos_} for t in doc]
    verbs = [t.text for t in doc if t.pos_ == "VERB"]
    nouns = [t.text for t in doc if t.pos_ == "NOUN"]

    return {
        "tokens": tokens,
        "verbs": verbs,
        "nouns": nouns
    }
\end{lstlisting}

% ===============================
% Text-Visualisierung (Verben markieren)
% ===============================
\section{Textvisualisierung}

Zur Verbesserung der Interpretierbarkeit werden erkannte Verben im Text visuell hervorgehoben.

\begin{lstlisting}[caption={nlp_analysis.py – Markierung von Verben im Text}]
def highlight_verbs(text):
    doc = nlp(text)
    out = []

    for token in doc:
        if token.pos_ == "VERB":
            out.append(f"<mark>{token.text}</mark>")
        else:
            out.append(token.text)

    return " ".join(out)
\end{lstlisting}

% ===============================
% Bildgenerierung
% ===============================
\section{Visuelle Ausgabe (Bildsuche)}

Die extrahierten Nomen werden als semantische Suchbegriffe verwendet, um passende Bilder aus einer externen Bildquelle abzurufen.

\begin{lstlisting}[caption={image_generator.py – Generierung semantischer Bildabfragen}]
import requests

def get_image_from_nouns(query):
    url = f"https://source.unsplash.com/800x600/?{query}"
    return url
\end{lstlisting}

% ===============================
% Streamlit Anwendung (Pipeline-Orchestrierung)
% ===============================
\section{Streamlit Anwendung}

Dieses Listing zeigt die vollständige Orchestrierung der KI-Pipeline innerhalb einer interaktiven Webanwendung.

\begin{lstlisting}[caption={app.py – Gesamtpipeline des multimodalen Systems}]
import streamlit as st
from speech_to_text import transcribe_audio
from nlp_analysis import analyze, highlight_verbs
from image_generator import get_image_from_nouns

st.title("Multimodaler Sprach-Assistent")

mode = st.radio("Eingabemodus", ["Text", "Audio"])
text = ""

if mode == "Text":
    text = st.text_input("Eingabe")
else:
    audio = st.file_uploader("Audio hochladen")
    if audio:
        text = transcribe_audio(audio)

if text:
    st.subheader("Linguistische Analyse")
    result = analyze(text)
    st.json(result)

    st.subheader("Verben hervorgehoben")
    st.markdown(highlight_verbs(text), unsafe_allow_html=True)

    if result["nouns"]:
        query = " ".join(result["nouns"][:3])
        st.image(get_image_from_nouns(query))
\end{lstlisting}